\documentclass[a4paper,10pt,lithuanian]{article}

\usepackage{lnfkstyle}

%\usepackage{graphicx}
\usepackage{amsmath}
\usepackage{amssymb}
\usepackage{url}

% extended math fonts
%\usepackage{yhmath}

\usepackage{srctex} % MikTex,Inverse Search in tex files from Yap screen

\newcommand{\ii}{\mathrm{i}} %imaginary unit
\newcommand{\ee}{\mathrm{e}} %exponent
\newcommand{\dd}{\mathrm{d}} %differential
\newcommand{\cl}[2]{\ensuremath{\mathit{Cl}_{#1,#2}}} % clifford algebra notation


%\DeclareFontFamily{U}{mathx}{\hyphenchar\font45}
%\DeclareFontShape{U}{mathx}{m}{n}{
%      <5> <6> <7> <8> <9> <10>
%      <10.95> <12> <14.4> <17.28> <20.74> <24.88>
%      mathx10
%      }{}
%\DeclareSymbolFont{mathx}{U}{mathx}{m}{n}
%\DeclareFontSubstitution{U}{mathx}{m}{n}
%\DeclareMathAccent{\widecheck}{0}{mathx}{"71}
%\DeclareMathAccent{\wideparen}{0}{mathx}{"75}


\newcommand{\reverse}[1]{\widetilde{#1}}
\newcommand{\gradeinverse}[1]{\wideparen{#1}}
%\newcommand{\cliffordconjugate}[1]{\wideparen{\widetilde{#1}}}

\def\d{\!\cdot\!} %narrow-space dot product
\def\w{\!\wedge\!} %narrow-space wedge product
\def\D{\cdot} %wide-space dot product
\def\W{\wedge} %wide-space wedge product


\newcommand{\e}[1]{\mathbf{e}_{#1}} % Clifford basis vectors in Cl(3,1)
\newcommand{\Ie}[1]{I\negthinspace\mathbf{e}_{#1}}%bivector in Cl(3,1)
\newcommand{\inv}[1]{\overline{#1}} %space inversion
\newcommand{\ba}{\mathbf{a}} %bold a
\newcommand{\bk}{\mathbf{k}}
\newcommand{\bx}{\mathbf{x}}
\newcommand{\bbf}{\mathbf{f}} %bold f

\def\m#1{\mathsf{#1}} % general multivector

\def\sR{\mathsf{R}} %sans serif Rotation
\def\sH{\mathsf{H}} %sans serif Hamiltonian
\def\sF{\mathsf{F}} %sans serif F
\def\sG{\mathsf{G}} %sans serif G
\def\cB{\mathcal{B}} %bivector




\begin{document}

\titlelt{Atvirkštiniai multivektoriai šešių dimensijų geometrinėse algebrose}

\titleen{Inverse multivectors in a six dimensional geometric algebras}


\author{\uline{Artūras Acus}$^{1}$, {Adolfas Dargys}$^{2}$}

\keywords{Geometrinė (Cliffordo) algebra,
atvirkštinis multivektorius, kompiuterinės algebros sistema}

\maketitle

\address{$^{1}$Teorinės fizikos ir astronomijos institutas,  Vilniaus universitetas, Saulėtekio 3, LT-10257 Vilnius}
\address{$^{2}$Nacionalinis fizinių ir technologijos mokslų centras,  Puslaidininkių fizikos institutas, Saulėtekio 3, LT-10257 Vilnius}

%\address{$^{2}$Fizinių ir technologijos mokslų centras, Savanorių pr. 231, 02300 Vilnius}

\rightaddress{arturas.acus@tfai.vu.lt}

\begin{multicols*}{2}

Geometrinė algebra (matematikų vadinama Cliffordo algebra)
apibendrina gerai žinomą vektorinį skaičiavimą, kuris plačiai naudojamas fizikoje. Jei vektorinis skaičiavimas tinka tik trimatėms erdvėms, tai įvedus multivektorius su geometrine algebra skaičiavimus galima atlikti bet kokios dimensijos ir signatūros erdvėse, tarp jų ir reliatyvistiame erdvėlaikyje. Svarbus geometrinės algebros elementas yra atvirkštinis multivektorius $A^{-1}$, kurio sandauga
iš multivektoriaus $A$ duoda vienetą.

Bendros formos multivektoriaus atvirkštinio radimas geometrinėje
algebroje tebėra iki galo neišspręsta ir labai aktuali problema
tiek matematikoje, tiek fizikoje. Kadangi geometrinės algebros
elementus visada galima pavaizduoti matricomis, tai iš principo
atvirkštinį multivektorių įmanoma rasti perėjus į matricinį
vaizdavimą. Tačiau toks kelias nėra natūralus, nes neišnaudoja
geometrinei algebrai būdingos rangų struktūros, kuri yra
prarandama perėjus į matricinį vaizdavimą. Bendras
atvirkštinių multivektorių formules galima būtų užrašyti
išskaidžius multivektorių ortogonalioje bazėje. Toks
metodas~\cite{key-1} turi keletą didelių trūkumų:
1)~išraiškos priklauso nuo pasirinktos bazės, 2)~jų
sudėtingumas, didėjant algebros dimensijai, auga
eksponentiškai, 3)~tos pačios dimensijos $n=p+q$, bet skirtingų
signatūrų $(p,q)$ geometrinėms algebroms $Cl_{p,q}$ formulės
skiriasi.

Tik gerokai vėliau buvo pastebėta~\cite{key-2}, kad algebroms,
kurių dimensija neviršija penkių,  $p+q\le 5$, atvirkštinius
multivektorius galima elegantiškai užrašyti bekordinatėje
formoje kaip pradinio mutivektoriaus $A$ ir įvairių jo rango
inversijų $A_{\bar r}$ (brūkšnelis virš indekso žymi, kad koeficiento prie
ranko $r$ bazinių elementų ženklas yra pakeičiamas į priešingą), tam tikrą sandaugą.
Pavyzdžiui,  algebroms kurių $n=3$,   ši formulė atrodo taip:
$A^{-1}=\frac{C(AC)_{\bar{3}}}{AC\,(AC)_{\bar{3}}}$, kur
$C=A_{\bar{1},\bar{2}}$. Deja, $n=6$ algebroms tokio pavidalo
formulės iki šiol niekam nepavyko rasti.

Pranešime mes parodome, kad atvirkštinių multivektorių
formules visgi galima užrašyti ir didesnės dimensijos
algebroms, jei paprastą multivektorių sandaugą
pakeisime  multivektorių multitiesine kombinacija. Tokią
kombinaciją mes išreikštai apskaičiavome $n=6$ geometrinei
algebrai~\cite{key-3}. Kadangi į gautą formulę įeina tik geometrinė sandauga
ir rangų inversija (involiucijos operacija), tai ta pati formulė
tinka bet kokios signatūros algebroms. Panašiai kaip ir $n=4$
atveju tokių formulių galima surasti ne vieną. Nors visos jos
leidžia apskaičiuoti atvirkštinius, tačiau skaičiavimo
efektyvumo požiūriu gerokai skiriasi. Mažiau rango
inversijų turinti formulė nebūtinai yra pati efektyviausia.
Daugiau involiucijų turinti formulė gali būti
spartesnė, nes kiekviename sumos naryje sumažėja galutiniame rezultate išsiprastinančių narių.

Atvirkštinių multivektorių išvedimas yra gan sudėtingas,
todėl sunkiai įsivaizduojamas be simbolinės kompiuterinės
algebros. Savo darbui mes sukūrėme specialiai tam skirtą {\it
Mathematica} sistemos paketą~\cite{key-4}, kuris leidžia greitai
ir patogiai sudauginti bet kurios geometrinės algebros
multivektorius, o taip pat rasti jų matricinius vaizdavimus.
Gautos formulės, kurios yra pateikiamos 1~lentelėje, buvo patikrintos simboliškai ir skaitiškai, apskaičiuojant įvairių algebrų atvirkštinius multivektorius. Išreikštinės formulės taip pat yra labai naudingos nustatant atvirkštinio multivektoriaus egzistavimo sąlygas.

\vspace{10pt} \noindent 1 lentelė. {\footnotesize Bendriausios
formos  multivektorių atvirkštinių formulių išraiškos geometrinėse
algebrose, kurių dimensija $p+q\le 6$. Kai $p+q=6$ antroji,
daugiau involiucijų turinti formulė, yra spartesnė.}
\begin{center}
\[\begin{array}{ll}
\cl{p}{q}&\quad  A^{-1} \\[5pt]\hline\\[-5pt]
\scriptstyle p+q=0&\quad\frac{B}{AB},\qquad B=1\\[4pt]
%
\scriptstyle p+q=1&\quad
\frac{B(AB)_{\bar{1}}}{AB\,(AB)_{\bar{1}}},\qquad B=1\\[4pt]
%
\scriptstyle  p+q=2&\quad \frac{C}{AC}
,\quad C=A_{\bar{1},\bar{2}}\\[4pt]
%
\scriptstyle p+q=3 &\quad
 \frac{C(AC)_{\bar{3}}}{AC\,(AC)_{\bar{3}}}
,\quad C=A_{\bar{1},\bar{2}}\\[4pt]
%
\scriptstyle p+q=4&\quad\frac{D}{AD},\qquad
D=A_{\bar{2},\bar{3}}
(AA_{\bar{2},\bar{3}})_{\bar{1},\bar{4}}\quad \\
   &\quad\textrm{\small arba}\qquad D=A_{\bar{1},\bar{2}}(A
A_{\bar{1},\bar{2}})_{\bar{3},\bar{4}}  \\[4pt]
%
\scriptstyle p+q=5&\quad\frac{D(AD)_{\bar{5}}}{AD\,
(AD)_{\bar{5}}},\qquad
D=A_{\bar{2},\bar{3}}
(AA_{\bar{2},\bar{3}})_{\bar{1},\bar{4}}\\[4pt]
%
\scriptstyle p+q=6&\quad\frac{G}{AG},\\
&
\!\!G=\frac{1}{3} A_{\bar{2},\bar{3},\bar{6}} \Bigl(H (H H)_{\bar{1},\bar{4},\bar{5}} + 2 \bigl(H_{\bar{4}} (H_{\bar{4}} H_{\bar{4}})_{\bar{1},\bar{4},\bar{5}}\bigr)_{\bar{4}}\Bigr)\\
&\textrm{\small arba}\\
&\!\!G=\frac{1}{3} A_{\bar{2},\bar{3},\bar{6}} \Bigl(\bigl(H (H H_{\bar{1},\bar{5}})_{\bar{4}}\bigr)_{\bar{1},\bar{5}}\\
&\hphantom{G=\frac{1}{3} A_{\bar{2},\bar{3},\bar{6}} \Bigl(\bigl(H (H H}
+ 2 \bigl(H_{\bar{4},\bar{5}} (H_{\bar{4},\bar{5}} H_{\bar{1},\bar{4}})_{\bar{4}}\bigr)_{\bar{1},\bar{4}}\Bigr)\\
&\quad\textrm{\small kur}\quad H=A A_{\bar{2},\bar{3},\bar{6}}
\end{array}
\]
\end{center}


\begin{thebibliography}{References}

\bibitem{key-1}
J.~P. Fletcher,
\newblock Clifford numbers and their inverses calculated using the matrix
  representation,
\newblock in {\em Applications of Geometric Algebra in Computer Science and
  Engineering}, edited by L.~Dorst, C.~Doran, and J.~Lasenby, pages 169--178,
  Birkh{\"a}user Boston, 2002.

\bibitem{key-2}
P.~{Dadbeh},
\newblock {Inverse and determinant in 0 to 5 dimensional {C}lifford algebra},
\newblock arXiv: 1104.0067  (Mar. 2011).

\bibitem{key-3}
A.~Acus and A.~Dargys,
\newblock submitted to Advances in Applied Clifford algebras

\bibitem{key-4}
A.~Acus and A.~Dargys,
\newblock {\it Mathematica} package, 2017,\newline
 \url{https://github.com/ArturasAcus/GeometricAlgebra}.

\end{thebibliography}



\end{multicols*}
\end{document}
